\section{The Constructible Side and the FLTZ Skeleton}

Throughout this section, let \(\Sigma\) be a rational polyhedral fan in
\(N_{\mathbb R}\). No smoothness assumption is imposed here, since the
constructions in this section depend only on the polyhedral data of
\(\Sigma\). The constructible side is not formulated on the toric variety
\(X_\Sigma\), but on the real vector space \(M_{\mathbb R}\). Under the natural
identification
\[
T^*M_{\mathbb R}
\cong
M_{\mathbb R}\times (M_{\mathbb R})^*
\cong
M_{\mathbb R}\times N_{\mathbb R},
\]
the fan \(\Sigma\subset N_{\mathbb R}\) appears as a source of prescribed
cotangent directions. Thus this section constructs the microlocal object that
controls the constructible side of the correspondence.

\subsection{The FLTZ Skeleton and Prescribed Singular Support}

We first define the FLTZ skeleton associated to \(\Sigma\). In the above
identification of the cotangent bundle, the first factor is the base space and
the second factor records cotangent directions. Thus the fan in
\(N_{\mathbb R}\) imposes a microlocal condition on sheaves on
\(M_{\mathbb R}\).

\begin{definition}[FLTZ skeleton]
For a cone \(\sigma\in\Sigma\), set
\[
\sigma^\perp
=
\{m\in M_{\mathbb R}\mid \langle m,n\rangle=0
\text{ for all }n\in\sigma\},
\qquad
M_\sigma=M/(\sigma^\perp\cap M).
\]
For \(\chi\in M_\sigma\), let \(\chi+\sigma^\perp\) denote the affine subspace
of \(M_{\mathbb R}\) determined by any lift of \(\chi\) to \(M\). This is
well-defined because two lifts differ by an element of \(\sigma^\perp\cap M\).

The FLTZ skeleton associated to \(\Sigma\) is the subset
\[
\Lambda_\Sigma
=
\bigcup_{\sigma\in\Sigma}
\bigcup_{\chi\in M_\sigma}
(\chi+\sigma^\perp)\times(-\sigma)
\subset T^*M_{\mathbb R}.
\]
\end{definition}

Equivalently,
\[
\Lambda_\Sigma
=
\bigcup_{\sigma\in\Sigma}
(\sigma^\perp+M)\times(-\sigma).
\]
Here \(\chi+\sigma^\perp\) records the base locus in \(M_{\mathbb R}\), while
\(-\sigma\) records the allowed cotangent directions.

\begin{proposition}\label{prop:fltz-conic}
The subset \(\Lambda_\Sigma\) is conic with respect to the cotangent fibers.
\end{proposition}

\begin{proof}
The conic action on \(T^*M_{\mathbb R}\cong M_{\mathbb R}\times N_{\mathbb R}\)
scales only the second factor. If
\[
(m,n)\in(\chi+\sigma^\perp)\times(-\sigma),
\]
then for every \(t\in\mathbb R_{\ge 0}\), the base point \(m\) is unchanged and
\(tn\in-\sigma\), since \(-\sigma\) is a cone. Hence
\[
(m,tn)\in(\chi+\sigma^\perp)\times(-\sigma)\subset\Lambda_\Sigma.
\]
Therefore \(\Lambda_\Sigma\) is conic in the cotangent direction.
\end{proof}

Using the notion of singular support recalled in Section~3, define
\[
Sh_{\Lambda_\Sigma}(M_{\mathbb R})
=
\{F\in Sh(M_{\mathbb R})\mid SS(F)\subset \Lambda_\Sigma\}.
\]
This is the category of constructible sheaves on \(M_{\mathbb R}\) whose
singular support is prescribed by the FLTZ skeleton.


\subsection{Singular Support of Polyhedral Sheaves}
We first record the singular-support estimate for costandard sheaves on
closed convex polyhedra.

\begin{lemma}[Singular support of a polyhedral costandard sheaf]
Let \(P\subset M_{\mathbb R}\) be a closed convex polyhedron, and let
\(i_P:P\hookrightarrow M_{\mathbb R}\) be the inclusion. Put
\[
\omega_P:=i_P^!k_{M_{\mathbb R}}.
\]
For a face \(G\preceq P\), define the normal cone
\[
N_GP:=
\{n\in N_{\mathbb R}\mid
\langle m-g,n\rangle\ge 0
\text{ for all }m\in P,\ g\in G\}.
\]
Then
\[
SS((i_P)_*\omega_P)
\subset
\bigcup_{G\preceq P}G\times(-N_GP)
\subset T^*M_{\mathbb R}.
\]
\end{lemma}

\begin{proof}
First consider \(k_P\). Stratify \(P\) by its relative open faces. Near a
point of \(\operatorname{relint}(G)\), the pair
\((M_{\mathbb R},P)\) is microlocally modeled by the tangent cone to \(P\)
along \(G\). The local computation for a closed convex polyhedral cone says
that the singular covectors of the constant sheaf on this cone are contained
in its polar cone. Hence, along the face \(G\), the singular covectors of
\(k_P\) are contained in \(N_GP\). Therefore
\[
SS(k_P)\subset
\bigcup_{G\preceq P}G\times N_GP .
\]

By Verdier duality,
\[
(i_P)_*\omega_P\simeq \mathbb D(k_P).
\]
Moreover,
\[
SS(\mathbb D F)=SS(F)^a,\qquad (x,\xi)^a=(x,-\xi).
\]
Thus
\[
SS((i_P)_*\omega_P)
=
SS(\mathbb D k_P)
\subset
\bigcup_{G\preceq P}G\times(-N_GP).
\]
\end{proof}




\begin{proposition}\label{prop:theta-singular-support}
Let \(\sigma\in\Sigma\), \(\chi\in M_\sigma\), and set
\(P_{\sigma,\chi}:=\chi+\sigma^\vee\). Then
\[
SS((i_{\sigma,\chi})_*\omega_{P_{\sigma,\chi}})
\subset
\bigcup_{\tau\preceq\sigma}(\chi+\tau^\perp)\times(-\tau)
\subset \Lambda_\Sigma .
\]
\end{proposition}

\begin{proof}
The faces of \(P_{\sigma,\chi}\) are indexed by faces
\(\tau\preceq\sigma\). The face corresponding to \(\tau\) is contained in
\(\chi+\tau^\perp\), and its normal cone is \(\tau\). Applying the
polyhedral singular-support estimate gives
\[
SS((i_{\sigma,\chi})_*\omega_{P_{\sigma,\chi}})
\subset
\bigcup_{\tau\preceq\sigma}(\chi+\tau^\perp)\times(-\tau).
\]
Since every face \(\tau\preceq\sigma\) is again a cone of \(\Sigma\), the
right-hand side is contained in \(\Lambda_\Sigma\).
\end{proof}

\begin{lemma}[Polyhedral microlocal dévissage]\label{lem:polyhedral-devissage}
Let \(\mathcal C_\Sigma\) be the stable idempotent-complete subcategory of
\(\operatorname{Sh}_{\Lambda_\Sigma}(M_{\mathbb R})\) generated by the
costandard sheaves
\[
(i_{\sigma,\chi})_*\omega_{\chi+\sigma^\vee},
\qquad \sigma\in\Sigma,\ \chi\in M_\sigma .
\]
Then every object of
\(\operatorname{Sh}_{\Lambda_\Sigma}(M_{\mathbb R})\) belongs to
\(\mathcal C_\Sigma\).
\end{lemma}

\begin{proof}[Sketch of proof]
The condition \(SS(F)\subset \Lambda_\Sigma\) implies that the singularities
of \(F\) can occur only along affine subspaces of the form
\(\chi+\sigma^\perp\), with conormal directions contained in \(-\sigma\).
Choose a polyhedral stratification of \(M_{\mathbb R}\) compatible with these
affine subspaces. By the constructible dévissage theorem, \(F\) is generated
by standard or costandard sheaves on the strata. The microlocal condition
excludes all strata and conormal directions not appearing in
\(\Lambda_\Sigma\). Therefore the remaining pieces are precisely
polyhedral costandard sheaves compatible with the fan \(\Sigma\), and these
are generated by the sheaves
\((i_{\sigma,\chi})_*\omega_{\chi+\sigma^\vee}\).
\end{proof}


\subsection{Theta Sheaves and Generation}

We now introduce the basic constructible objects attached to the cones of
\(\Sigma\). For \(\sigma\in\Sigma\), let
\(\sigma^\vee=\{m\in M_{\mathbb R}\mid \langle m,n\rangle\ge 0
\text{ for all }n\in\sigma\}\). For \(\chi\in M_\sigma\), choose a lift of
\(\chi\) to \(M\) and set
\(P_{\sigma,\chi}:=\chi+\sigma^\vee\subset M_{\mathbb R}\). This subset is
independent of the chosen lift, since \(\sigma^\perp\) is the lineality space
of \(\sigma^\vee\). Let
\(i_{\sigma,\chi}:P_{\sigma,\chi}\hookrightarrow M_{\mathbb R}\) be the
inclusion. The constructible theta sheaf associated to \((\sigma,\chi)\) is
\[
\Theta(\sigma,\chi):=(i_{\sigma,\chi})_*\omega_{P_{\sigma,\chi}}.
\]
Thus \(\Theta(\sigma,\chi)\) is the costandard sheaf supported on the
translated dual cone \(P_{\sigma,\chi}\). By \cref{prop:theta-singular-support}, one has
\(SS(\Theta(\sigma,\chi))\subset \Lambda_\Sigma\), and hence
\(\Theta(\sigma,\chi)\in
\operatorname{Sh}_{\Lambda_\Sigma}(M_{\mathbb R})\).

\begin{theorem}[FLTZ generation]\label{thm:fltz-generation}
The category \(\operatorname{Sh}_{\Lambda_\Sigma}(M_{\mathbb R})\) is generated
by the theta sheaves:
\[
\operatorname{Sh}_{\Lambda_\Sigma}(M_{\mathbb R})
=
\left\langle
\Theta(\sigma,\chi)\mid \sigma\in\Sigma,\ \chi\in M_\sigma
\right\rangle .
\]
Here the right-hand side denotes the smallest stable idempotent-complete
subcategory containing all \(\Theta(\sigma,\chi)\).
\end{theorem}

\begin{proof}
Let \(\mathcal C_\Sigma\) be the stable idempotent-complete subcategory
generated by the theta sheaves. Since
\(SS(\Theta(\sigma,\chi))\subset \Lambda_\Sigma\) for every
\(\sigma\in\Sigma\) and \(\chi\in M_\sigma\), and since
\(\operatorname{Sh}_{\Lambda_\Sigma}(M_{\mathbb R})\) is closed under shifts,
cones, finite direct sums, and retracts, we have
\(\mathcal C_\Sigma\subset
\operatorname{Sh}_{\Lambda_\Sigma}(M_{\mathbb R})\).

Conversely, \cref{lem:polyhedral-devissage} says that every object of
\(\operatorname{Sh}_{\Lambda_\Sigma}(M_{\mathbb R})\) belongs to the stable
idempotent-complete subcategory generated by the sheaves
\((i_{\sigma,\chi})_*\omega_{\chi+\sigma^\vee}
=\Theta(\sigma,\chi)\). Hence
\(\operatorname{Sh}_{\Lambda_\Sigma}(M_{\mathbb R})\subset
\mathcal C_\Sigma\). The two inclusions give the desired equality.
\end{proof}